\documentclass[11pt]{article}
\usepackage{amsmath, amssymb, bm}
\usepackage{geometry}
\usepackage{hyperref}
\usepackage{parskip} 
\usepackage{booktabs}
\usepackage{array}

\geometry{margin=1.1in}

\title{\textbf{RoboRacer Control: PID, MPPI, and Nonlinear MPC}\\[0.4em]
\large Mathematical Overview}
\author{}
\date{}

\begin{document}
\maketitle
\tableofcontents
\newpage

% ─────────────────────────────────────────────────────────────────────────────
\section{Vehicle Model}

All three controllers share the same underlying vehicle model: the
\textbf{kinematic bicycle model}.  It treats the car as a single-track rigid
body with no tire slip, which is valid at low-to-moderate speeds where lateral
acceleration is small enough that the tires remain near the linear region of
their force--slip curves.

\subsection{States and Inputs}

\[
\bm{x} = \begin{pmatrix} X \\ Y \\ \psi \\ v \end{pmatrix}, \qquad
\bm{u} = \begin{pmatrix} \delta \\ a \end{pmatrix}
\]

\begin{center}
\begin{tabular}{lll}
\toprule
Symbol & Description & Unit \\
\midrule
$X, Y$ & Global position & m \\
$\psi$ & Heading angle (yaw) & rad \\
$v$    & Longitudinal speed & m s$^{-1}$ \\
$\delta$ & Front-wheel steering angle & rad \\
$a$    & Longitudinal acceleration & m s$^{-2}$ \\
$L$    & Wheelbase & m \\
\bottomrule
\end{tabular}
\end{center}

\subsection{Continuous-Time Dynamics}

\begin{align}
\dot{X}   &= v \cos\psi \label{eq:xdot}\\
\dot{Y}   &= v \sin\psi \label{eq:ydot}\\
\dot{\psi} &= \frac{v}{L}\tan\delta \label{eq:psidot}\\
\dot{v}   &= a \label{eq:vdot}
\end{align}

Equation~\eqref{eq:psidot} is derived from Ackermann geometry: for a bicycle
model the instantaneous turning radius is $R = L / \tan\delta$, so the yaw rate
is $\dot{\psi} = v/R$.

\subsection{Constraints}

\[
\delta \in [-\delta_{\max},\; \delta_{\max}], \qquad
a     \in [a_{\min},\; a_{\max}], \qquad
v     \in [0,\; v_{\max}]
\]

with $\delta_{\max} = 0.44\,\text{rad}$, $a_{\min} = -6.15$, $a_{\max} = 6.15\,\text{m\,s}^{-2}$,
$v_{\max} = 15\,\text{m\,s}^{-1}$, and $L = 0.324\,\text{m}$.


% ─────────────────────────────────────────────────────────────────────────────
\section{Raceline and Reference Trajectory}

A pre-computed raceline CSV provides the reference:
\[
\mathcal{R} = \{(X_i,\, Y_i,\, \psi_i,\, v_{\text{ref},i},\, s_i)\}_{i=0}^{M-1}
\]
where $s_i$ is the cumulative arc length at waypoint $i$.

\subsection{Windowed Closest-Point Search}

The \textbf{closest waypoint index} $i^*$ is the nearest neighbour of the car's
position $(X, Y)$ in the raceline.  A pure exhaustive search is correct but
fails at \emph{self-intersections} of the raceline: on the figure-8 track the
loops cross at a single physical point that maps to two different arc lengths,
and an exhaustive search can flip between branches when the car is near the
crossing.

To prevent this, after the first call the search is restricted to a window of
$\pm 100$ indices around the previous $i^*$.  This forces monotonic progress
along the arc and eliminates the crossing ambiguity at the cost of requiring
that the car never moves more than one window's worth of arc length between
control steps (at $\Delta t = 25$\,ms and 100 indices $\approx 2.9$\,m of arc,
this corresponds to speeds well above $v_{\max}$).

\subsection{Lookahead}

Lookahead reference points at horizon step $k$ are obtained by projecting $k$
steps forward in arc length:
\[
s_{\text{ref}}(k) = s_{i^*} + k \cdot v_{\text{look}} \cdot \Delta t,
\qquad
v_{\text{look}} = \max(\hat{v},\; v_{\text{look,min}})
\]
where $\hat{v}$ is the current (estimated) car speed and
$v_{\text{look,min}} = 0.8$\,m\,s$^{-1}$ is a floor that keeps the lookahead
non-degenerate when stationary.  Looking the reference up at the
\emph{car's actual speed} rather than the \emph{reference speed}
$v_{\text{ref},i^*}$ keeps the reference spread short during startup
acceleration; if the slow car were given a reference that already extends
several metres ahead, the controller would chase a point it cannot
physically reach in the horizon, biasing the solution toward unrealistic
inputs.  Once the car is up to speed the two choices coincide.


% ─────────────────────────────────────────────────────────────────────────────
\section{MPPI --- Model Predictive Path Integral Control}

\subsection{Conceptual Overview}

MPPI (Williams et al., 2017) is a \textbf{sampling-based} MPC variant rooted in
stochastic optimal control and information theory.  Instead of solving an
optimisation problem, it runs $K$ random \emph{rollouts} of the vehicle model
forward in time, scores each one with a cost function, and then computes the
optimal control as a cost-weighted average of all the rollout inputs.

On the RoboRacer, at every 25\,ms control step:
\begin{enumerate}
  \item $K = 300$ random perturbation sequences are sampled.
  \item Each is propagated through the bicycle model for $N = 20$ steps.
  \item Each trajectory is scored: trajectories that stay close to the raceline
        receive high importance weight; trajectories that deviate are down-weighted.
  \item The optimal input is a soft mixture of all sequences, dominated by the
        best-performing ones.
  \item Only the first input of the mixture is applied to the car; the full
        sequence is shifted and reused as the warm start for the next step.
\end{enumerate}

\subsection{Discrete-Time Dynamics (Euler)}

MPPI propagates an ensemble of states simultaneously using vectorised forward
Euler integration:

\begin{align}
X_{t+1}^{(k)}   &= X_t^{(k)} + v_t^{(k)}\cos\psi_t^{(k)}\,\Delta t \\
Y_{t+1}^{(k)}   &= Y_t^{(k)} + v_t^{(k)}\sin\psi_t^{(k)}\,\Delta t \\
\psi_{t+1}^{(k)} &= \psi_t^{(k)} + \frac{v_t^{(k)}}{L}\tan\delta_t^{(k)}\,\Delta t \\
v_{t+1}^{(k)}   &= \mathrm{clip}\!\left(v_t^{(k)} + a_t^{(k)}\,\Delta t,\; 0,\; v_{\max}\right)
\end{align}

The $\mathrm{clip}$ in the velocity update is a practical bound: it is
non-differentiable, but MPPI never requires gradients through the dynamics.

\subsection{Algorithm}

Let $\bm{U} = [\bm{u}_0, \dots, \bm{u}_{N-1}] \in \mathbb{R}^{N \times 2}$
be the current nominal (warm-start) sequence and $\lambda > 0$ the temperature.

\paragraph{Step 1 -- Sample perturbations.}
\[
\bm{\varepsilon}^{(k)} \sim \mathcal{N}(\bm{0},\, \bm{\Sigma}), \qquad k = 1,\dots,K
\]
with $\bm{\Sigma} = \mathrm{diag}(\sigma_\delta^2, \sigma_a^2)$,
$\sigma_\delta = 0.3\,\text{rad}$, $\sigma_a = 1.5\,\text{m\,s}^{-2}$.  Note that
$\sigma_\delta$ is comparable to $\delta_{\max} = 0.44$\,rad: by 1$\sigma$ the
samples already span most of the steering range, ensuring the corner-radius
inputs needed on the figure-8 are present in the rollout cloud.

\paragraph{Step 2 -- Form and clip perturbed sequences.}
\[
\bm{V}^{(k)} = \mathrm{clip}\!\left(\bm{U} + \bm{\varepsilon}^{(k)},\; \bm{u}_{\min},\; \bm{u}_{\max}\right)
\]

\paragraph{Step 3 -- Roll out and accumulate cost.}
Starting from the current state $\bm{x}_0$, propagate each sequence and
accumulate the running cost:
\[
J^{(k)} = \sum_{t=0}^{N-1} c\!\left(\bm{x}_t^{(k)},\; \bm{V}_t^{(k)}\right)
\]

\paragraph{Step 4 -- Compute importance weights.}
Subtract the minimum for numerical stability, then normalise:
\[
\beta = \min_k J^{(k)}, \qquad
\tilde{w}_k = \exp\!\left(-\frac{J^{(k)} - \beta}{\lambda}\right), \qquad
w_k = \frac{\tilde{w}_k}{\sum_j \tilde{w}_j}
\]
Low temperature $\lambda$ concentrates weight on the best rollouts (greedy);
high temperature spreads it more evenly (exploratory).

\paragraph{Step 5 -- Update nominal sequence.}
\[
\bm{U}^* = \sum_{k=1}^{K} w_k \bm{V}^{(k)}
\]

\paragraph{Step 6 -- Apply and shift.}
Apply $\bm{u}_0^*$ to the car.  Shift the sequence left by one step and zero-pad
the tail to form the warm start for the next control iteration.

\subsection{Cost Function}

\[
c(\bm{x}_t, \bm{u}_t) =
  W_{\text{CTE}} \underbrace{\left\|\bm{p}_t - \bm{p}_{\text{ref}}(t)\right\|_2}_{\text{cross-track error}}
+ W_h \underbrace{\left|\psi_t - \psi_{\text{ref}}(t)\right|_{\pi}}_{\text{heading error}}
+ W_v \underbrace{\left|v_t - v_{\text{ref}}(t)\right|}_{\text{speed error}}
+ W_s\, \underbrace{(\delta_t - \delta_{t-1})^2}_{\text{steering rate}}
\]

where $|\cdot|_\pi$ denotes the wrapped absolute angle difference and
$\bm{p}_{\text{ref}}(t)$, $\psi_{\text{ref}}(t)$, $v_{\text{ref}}(t)$ are
looked up from the raceline at the projected arc position.

A subtle but important detail is that the steering term penalises the
\emph{rate} of change $(\delta_t - \delta_{t-1})^2$ rather than the absolute
magnitude $\delta_t^2$.  Penalising the magnitude would discourage holding
the sustained $\delta \approx 0.21$\,rad required to follow the figure-8 loops
(radius $R \approx 1.5$\,m, $\delta_{\text{steady}} = \arctan(L/R)$), causing
systematic understeer.  Penalising the rate suppresses bang-bang jitter
without fighting steady cornering inputs.  The previous control $\delta_{t-1}$
is seeded from the actual last applied input at step $t=0$ rather than from
zero, anchoring the rate term to physical continuity.

\medskip
\begin{center}
\begin{tabular}{lll}
\toprule
Weight & Value & Meaning \\
\midrule
$W_{\text{CTE}}$ & 23.0   & Penalise lateral deviation from the raceline \\
$W_h$            & 20.0   & Penalise heading misalignment \\
$W_v$            & 0.5    & Penalise speed error relative to $v_{\text{ref}}$ \\
$W_s$            & 0.0005 & Penalise steering rate $(\delta_t - \delta_{t-1})^2$ \\
\bottomrule
\end{tabular}
\end{center}

Other implementation parameters:

\medskip
\begin{center}
\begin{tabular}{lll}
\toprule
Parameter & Value & Meaning \\
\midrule
$K$              & 300    & Number of rollouts per control step \\
$N$              & 20     & Horizon length (steps) \\
$\Delta t$       & 0.025\,s & Prediction step size (matches main loop period) \\
$\lambda$        & 20.0   & MPPI temperature \\
$\sigma_\delta$  & 0.3\,rad & Steering perturbation std \\
$\sigma_a$       & 1.5\,m\,s$^{-2}$ & Acceleration perturbation std \\
\bottomrule
\end{tabular}
\end{center}

\subsection{Warm-Start from Raceline Curvature}

On the very first control iteration the nominal sequence $\bm{U}$ has no prior
solve to inherit from.  Initialising it to zero biases all rollouts toward
straight-line motion, which catastrophically lags the controller through the
first corner.  Instead, the initial $\bm{U}$ is derived directly from the
raceline geometry.

For each horizon step $k$, two arc-length positions are sampled,
$s_k = s_{i^*} + k v_{\text{look}} \Delta t$ and
$s_{k+1} = s_{i^*} + (k+1) v_{\text{look}} \Delta t$, with corresponding
raceline headings $\psi_k$ and $\psi_{k+1}$.  Inverting the bicycle yaw-rate
equation~\eqref{eq:psidot},
\[
\dot{\psi} = \frac{v}{L}\tan\delta
\;\;\Longrightarrow\;\;
\delta_k = \mathrm{atan2}\!\left(L \cdot (\psi_{k+1} - \psi_k)_{\pi},\; v_{\text{look}}\,\Delta t\right)
\]
gives the steering input that would follow the raceline arc segment.  The
\texttt{atan2} form is used over $\arctan$ to avoid division-by-zero at small
$v_{\text{look}}$.  The acceleration component is left at zero; the speed
controller catches up within the first second.

\subsection{Velocity Tracking Change}

Earlier iterations of MPPI used $v_{\text{ref},i^*}$ to set the lookahead
spread.  At startup this caused the reference to leap several metres ahead of
a stationary car, producing rollouts that all tried to teleport rather than
accelerate.  The current code uses $v_{\text{look}} = \max(\hat{v}, 0.8)$,
which scales the reference spread with actual progress.  This and matching
$\Delta t = 0.025$\,s to the main loop period (it was previously 0.05\,s) were
the two largest single-parameter improvements during development.


% ─────────────────────────────────────────────────────────────────────────────
\section{Nonlinear MPC}

\subsection{Conceptual Overview}

Nonlinear MPC frames vehicle control as a \textbf{constrained optimisation
problem} solved at every timestep.  Rather than sampling random trajectories,
it \emph{solves} for the single best input sequence that minimises a quadratic
cost subject to the vehicle dynamics and hard input/state constraints.

On the RoboRacer, at every 25\,ms control step:
\begin{enumerate}
  \item The current Vicon pose is fed to the acados solver as the initial condition.
  \item The reference trajectory for the next $N$ prediction steps is looked up from
        the raceline.
  \item The solver performs one SQP-RTI step: it linearises the nonlinear dynamics
        around the warm-started trajectory, solves the resulting quadratic programme
        (QP) using HPIPM, and returns the optimal input.
  \item Only the first input of the sequence is applied; the rest warm-start the
        next solve.
\end{enumerate}

\subsection{Continuous-Time Dynamics}

Unlike MPPI, the MPC uses the continuous-time ODE \eqref{eq:xdot}--\eqref{eq:vdot}
directly. acados integrates it numerically using explicit Runge--Kutta 4 (ERK4):

\[
\bm{x}_{k+1} = \bm{x}_k + \frac{\Delta t}{6}
  \left(\bm{f}_1 + 2\bm{f}_2 + 2\bm{f}_3 + \bm{f}_4\right)
\]

where $\bm{f}_1, \dots, \bm{f}_4$ are the standard RK4 stage derivatives of
$\dot{\bm{x}} = f(\bm{x}, \bm{u})$.  This is more accurate than forward Euler
and does not require the velocity $\mathrm{clip}$; speed bounds are instead
expressed as explicit inequality constraints:
\[
0 \leq v_k \leq v_{\max} \qquad \forall\, k = 1, \dots, N
\]

\subsection{Optimal Control Problem}

At time $t$ with current state $\bm{x}_0$:

\begin{align}
\min_{\bm{u}_0,\,\dots,\,\bm{u}_{N-1}} \quad
  &\sum_{k=0}^{N-1} \ell(\bm{x}_k,\, \bm{u}_k)
  \;+\; \ell_N(\bm{x}_N) \label{eq:ocp}\\
\text{subject to} \quad
  &\bm{x}_{k+1} = F(\bm{x}_k, \bm{u}_k) \tag{dynamics}\\
  &\bm{u}_k \in [\bm{u}_{\min},\, \bm{u}_{\max}] \tag{input bounds}\\
  &v_k \in [0,\, v_{\max}] \tag{speed bound}\\
  &\bm{x}_0 = \bm{x}_{\text{current}} \tag{initial condition}
\end{align}

where $F(\bm{x}_k, \bm{u}_k)$ denotes one ERK4 step of the bicycle ODE.

\subsection{Cost Function}

A Linear Least-Squares (LINEAR\_LS) cost is used:
\[
\ell(\bm{x}_k, \bm{u}_k)
  = \left\|\bm{y}_k - \bm{y}_{\text{ref},k}\right\|_{\bm{W}}^2
  = \left(\bm{y}_k - \bm{y}_{\text{ref},k}\right)^{\!\top}
    \bm{W}
    \left(\bm{y}_k - \bm{y}_{\text{ref},k}\right)
\]

with the residual vector
\[
\bm{y}_k = \begin{pmatrix} X_k \\ Y_k \\ \psi_k \\ v_k \\ \delta_k \\ a_k \end{pmatrix},
\qquad
\bm{y}_{\text{ref},k} = \begin{pmatrix} X_{\text{ref},k} \\ Y_{\text{ref},k} \\ \psi_{\text{ref},k} \\ v_{\text{ref},k} \\ 0 \\ 0 \end{pmatrix},
\qquad
\bm{W} = \mathrm{diag}(W_{\text{CTE}},\, W_{\text{CTE}},\, W_h,\, W_v,\, W_s,\, W_a)
\]

The zero reference for $\delta$ and $a$ means the solver penalises control effort
(regularisation), keeping the trajectory smooth.  The terminal cost uses the
same structure with $\bm{W}_N = \mathrm{diag}(W_{\text{CTE}}, W_{\text{CTE}}, W_h, W_v)$.

Note that LINEAR\_LS penalises the \emph{absolute magnitude} $\delta_k^2$ rather
than the rate $(\delta_k - \delta_{k-1})^2$ used in MPPI.  Encoding a rate
penalty cleanly in LINEAR\_LS would require augmenting the state with the
previous control; in practice the RTI warm-start provides enough temporal
smoothness on its own that this has not been necessary.

\medskip
\begin{center}
\begin{tabular}{lll}
\toprule
Weight & Value & Meaning \\
\midrule
$W_{\text{CTE}}$ & 23.0 & Position tracking (applied to both $X$ and $Y$ residuals) \\
$W_h$            & 20.0 & Heading tracking \\
$W_v$            & 0.5  & Speed tracking \\
$W_s$            & 1.0  & Steering regularisation (absolute magnitude) \\
$W_a$            & 0.1  & Acceleration regularisation \\
\bottomrule
\end{tabular}
\end{center}

Other implementation parameters:

\medskip
\begin{center}
\begin{tabular}{lll}
\toprule
Parameter & Value & Meaning \\
\midrule
$N$        & 20       & Horizon length (steps) \\
$\Delta t$ & 0.025\,s & Prediction step (matches main loop period) \\
QP solver  & HPIPM PC & Partial-condensing interior point \\
NLP solver & SQP-RTI  & One linearisation per control step \\
Integrator & ERK4     & Explicit Runge--Kutta 4 \\
\bottomrule
\end{tabular}
\end{center}

\subsection{Reference Heading Unwrapping}

The continuous-time dynamics do not wrap $\psi$ --- a discontinuity at $\pm\pi$
would break the linearisation that SQP-RTI relies on.  The raceline's
$\psi_{\text{ref}}$ values, however, are stored in $[-\pi, \pi]$ from
\texttt{atan2}, and the figure-8 contains a $\pm\pi$ wrap roughly halfway
through the loop.  Feeding a wrapped reference into a non-wrapping state model
produces a huge residual at the wrap point.

The MPC therefore \textbf{unwraps} the reference incrementally: starting from
$\tilde{\psi}_0 = \psi_{\text{state}}$, each subsequent reference is built as
\[
\tilde{\psi}_{k+1} = \tilde{\psi}_k + \big(\psi_{\text{ref},k+1} - \tilde{\psi}_k\big)_\pi
\]
where $(\cdot)_\pi$ wraps an angle into $[-\pi, \pi]$.  This produces a
continuous reference $\tilde{\psi}_k$ in $\mathbb{R}$ that stays close to the
predicted $\psi_k$ trajectory throughout the horizon.

\subsection{Warm-Start from Raceline Curvature}

The same raceline-curvature initialisation used in MPPI is applied to the MPC
solver on its first iteration, before any previous solve exists to warm-start
from.  Stage inputs are seeded with
$\delta_k = \mathrm{atan2}(L\,\Delta\psi_k,\; v_{\text{look}}\Delta t)$ and
$a_k = 0$; stage states are seeded by walking along the raceline waypoints
with the unwrapped reference heading.  Acados then propagates from this guess
in subsequent solves via internal warm-starting.

\subsection{SQP Real-Time Iteration (SQP-RTI)}

Solving the full NLP \eqref{eq:ocp} to convergence at 40\,Hz is infeasible for
large horizons.  The \textbf{Real-Time Iteration} (Diehl et al., 2005) scheme
replaces the full NLP solve with a single SQP step per control cycle.

At each step, the NLP is linearised around the current warm-start trajectory
$(\bar{\bm{X}}, \bar{\bm{U}})$ to obtain a \textbf{quadratic programme}:

\begin{align}
\min_{\Delta\bm{U}} \quad
  &\frac{1}{2}\,\Delta\bm{U}^{\!\top} H\,\Delta\bm{U}
  + \bm{g}^{\!\top}\Delta\bm{U} \\
\text{subject to} \quad
  &C\,\Delta\bm{U} = \bm{c} \tag{linearised dynamics}\\
  &\bm{u}_{\min} - \bar{\bm{U}} \leq \Delta\bm{U} \leq \bm{u}_{\max} - \bar{\bm{U}} \tag{input bounds}
\end{align}

where $H$ is the Gauss--Newton Hessian approximation and $\bm{g}$ the gradient
of the cost.  The QP is solved by \textbf{HPIPM} (partial condensing), and the
solution $\Delta\bm{U}^*$ is added to the warm start.

Because the initial state changes only slightly between control steps (the car
has moved 25\,ms worth), the warm-started trajectory is already close to the new
optimum and one SQP step is sufficient to track it accurately.


% ─────────────────────────────────────────────────────────────────────────────
\section{PID Baseline (Stanley-style)}

\subsection{Conceptual Overview}

The PID controller is the classical baseline against which MPPI and MPC are
benchmarked.  Unlike the two predictive controllers, it has no internal model
and no horizon: it produces a steering command directly from the
\emph{current} cross-track and heading errors.  Structurally it is a
\emph{Stanley-style} steering law, augmented with an integral term:
\[
\delta = -\big(K_h \psi_e \;+\; K_e e \;+\; K_{\dot e}\, \dot e \;+\; K_{\int} \!\!\int e\, dt\big)
\]
The overall negative sign converts a positive (left-of-track) lateral error
into a right-steer command, given the sign convention of $\delta$ in the
bicycle model.  All four gains are kept positive; flipping the sign of any
single one would reverse the corresponding correction direction.

\subsection{Signed Cross-Track Error}

The plain Euclidean distance $\|\bm{p} - \bm{p}_{\text{ref}}\|$ used by MPPI
and MPC for cost is unsigned and cannot indicate which side of the track the
car is on.  PID needs a \textbf{signed} CTE to know which direction to
correct.  Decomposing the position offset along the raceline's local frame,
\[
\hat{\bm{t}}_{\text{ref}} = (\cos\psi_{\text{ref}},\, \sin\psi_{\text{ref}}),
\qquad
\hat{\bm{n}}_{\text{ref}} = (-\sin\psi_{\text{ref}},\, \cos\psi_{\text{ref}}),
\]
gives the signed CTE as the projection of the offset onto the left-normal:
\[
e = (\bm{p} - \bm{p}_{\text{ref}}) \cdot \hat{\bm{n}}_{\text{ref}}
  = -\sin\psi_{\text{ref}}\,(X - X_{\text{ref}}) + \cos\psi_{\text{ref}}\,(Y - Y_{\text{ref}})
\]
By convention $e > 0$ when the car is to the \emph{left} of the track.

\subsection{Derivative and Integral Terms}

The derivative term $\dot e$ is approximated by a finite difference between
successive control steps,
\[
\dot e \approx \frac{e_t - e_{t-1}}{\Delta t},
\]
with $\Delta t$ clamped to $[10^{-6}, 0.2]\,\text{s}$ to prevent first-step
spikes and stale-frame outliers from dominating the derivative.

The integral term $\int e\, dt$ is implemented as a discrete sum with
\textbf{clamping anti-windup}: the running integral is bounded so that the
integral contribution to $\delta$ cannot exceed a configurable cap
$I_{\max}$, and the integral is frozen entirely while the total steering
command is saturated at $\pm\delta_{\max}$.  The combination prevents the
integrator from accumulating during corner saturation and then overshooting
on exit.

\subsection{Heading Error Term and Stanley Analogue}

The Stanley controller in its original form is
\[
\delta = \psi_e + \arctan\!\left(\frac{k\, e}{v}\right).
\]
For small CTE the second term linearises to $k\,e/v$, motivating
$K_e \approx k/v$ as a starting gain.  At $v \approx 4.5\,\text{m/s}$ with
$k = 2$ this gives $K_e \approx 0.44$\,rad/m.  The heading term has gain
$K_h = 1.0$ (unit gain) so that small heading errors map directly to
proportional steering corrections; the closed-loop heading time constant of
this term is $\tau_h = L / (v K_h)$, which at $v = 4.5\,\text{m/s}$ evaluates
to $\tau_h \approx 72\,\text{ms}$.

\subsection{Parameters}

Default tuning (modifiable from the CLI):

\medskip
\begin{center}
\begin{tabular}{lll}
\toprule
Gain & Value & Meaning \\
\midrule
$K_h$        & 0.6  & Heading-error P gain (rad/rad) \\
$K_e$        & 0.6  & Cross-track P gain (rad/m) \\
$K_{\dot e}$ & 0.05 & Cross-track D gain (rad$\cdot$s/m) \\
$K_{\int}$   & 0.0  & Cross-track I gain (rad/(m$\cdot$s)) \\
$I_{\max}$   & 0.10\,rad & Anti-windup cap on integral contribution \\
\bottomrule
\end{tabular}
\end{center}

The integral gain is left at zero by default: the figure-8 is symmetric and
on a well-tuned PD the steady-state offset is negligible.  $K_{\int}$ is only
introduced when persistent one-sided drift is observed (e.g. from chassis
misalignment).

\subsection{Tuning Procedure}

Gains are tuned in a fixed order to isolate one closed-loop effect at a time:

\begin{enumerate}
  \item Set $K_e, K_{\dot e}, K_{\int}$ to zero; tune $K_h$ alone on a straight
        until heading error settles below $0.1$\,rad with no oscillation.
  \item Add $K_e$ to remove lateral drift; expect to revisit $K_h$ once $K_e$
        is active because the two terms partially compensate each other.
  \item Add $K_{\dot e}$ if the closed loop is now oscillatory but
        non-divergent.  D-gain is bounded above by measurement-noise
        amplification: roughly $K_{\dot e} < K_e / 2$ at $40$\,Hz.
  \item Add $K_{\int}$ only if a persistent one-sided offset remains.
\end{enumerate}


% ─────────────────────────────────────────────────────────────────────────────

\section{Throttle Control}

None of the three controllers directly commands a wheel torque.  All use a
\textbf{feedforward + proportional speed controller} that maps the raceline
reference velocity to a dimensionless throttle integer (0--200):

\[
\tau = \mathrm{clip}\!\left(K_{\mathrm{ff}}\, v_{\mathrm{ref}} + K_p\,(v_{\mathrm{ref}} - \hat{v}),\; 0,\; \tau_{\max}\right)
\]

where $\hat{v}$ is the speed estimate from Vicon position differences
(EMA-smoothed), $K_{\mathrm{ff}} = 20$ and $K_p = 5$ are default gains.
The MPPI acceleration output $a$ and the MPC acceleration command $a^*$ are
not used for throttle for four compounding reasons:

\begin{enumerate}
  \item \textbf{No calibrated actuator model.}  Both predictive controllers
        output $a$ in m\,s$^{-2}$, but the hardware input is a dimensionless
        throttle integer sent to an ESC.  The mapping throttle $\to$ torque
        $\to$ acceleration varies with speed, battery voltage, and surface
        friction.  Without empirical system identification a direct
        $a \to \tau$ conversion is unreliable.

  \item \textbf{Noisy velocity estimate.}  $\hat{v}$ is computed from finite
        differences of Vicon positions and is EMA-smoothed but still noisy.
        That noise enters the MPPI/MPC initial state, propagates through all
        rollouts or linearisations, and produces a noisy $a$ command.  A simple
        proportional controller on $v_{\mathrm{ref}} - \hat{v}$ is far more
        forgiving of this noise than closing the throttle loop through the
        optimiser.

  \item \textbf{Speed profile already solved offline.}  The raceline CSV encodes
        a pre-optimised speed profile (slow for tight corners, fast on
        straights).  Re-optimising speed online adds complexity without benefit
        when the reference profile is already well-designed; the controller
        need only track it.

  \item \textbf{Decoupled tuning.}  At low slip angles, lateral and longitudinal
        dynamics interact weakly.  Keeping the speed loop separate means
        $K_{\mathrm{ff}}$ and $K_p$ can be tuned independently of the MPPI
        hyper-parameters or MPC weights, making both easier to diagnose.
\end{enumerate}

\medskip

The steering servo command (integer 120--880) is derived from the controller's
steering angle output in all three cases:
\[
\tau_{\delta} = \mathrm{clip}\!\left(\tau_{\text{center}} - K_{\text{srv}}\,\delta + \tau_{\text{trim}},\;
  \tau_{\min},\; \tau_{\max}\right)
\]
with $\tau_{\text{center}} = 512$, $K_{\text{srv}} = 1200$\,servo-units/rad,
$\tau_{\text{trim}} = 20$.


% ─────────────────────────────────────────────────────────────────────────────
\section{Future Model Improvements}

\subsection{Dynamic Bicycle Model}

The kinematic model assumes zero tire slip, which is only valid when lateral
acceleration $a_y \ll \mu g$.  At higher speeds the tires develop meaningful
slip angles and the kinematic model underestimates the required steering.

The \textbf{dynamic bicycle model} adds lateral velocity $v_y$ and yaw rate
$r = \dot{\psi}$ as states, and replaces the kinematic steering constraint with
force balance equations:

\[
\bm{x} = \begin{pmatrix} X \\ Y \\ \psi \\ v_x \\ v_y \\ r \end{pmatrix},
\qquad
m\dot{v}_y = F_{y,f} + F_{y,r} - m v_x r, \qquad
I_z\dot{r} = L_f F_{y,f} - L_r F_{y,r}
\]

where $m$ is vehicle mass, $I_z$ yaw moment of inertia, $L_f$/$L_r$ are
distances from CG to front/rear axles, and $F_{y,f}$, $F_{y,r}$ are lateral
tire forces.  For a linear tire model: $F_{y,f} = -C_f\alpha_f$,
$F_{y,r} = -C_r\alpha_r$, with slip angles

\[
\alpha_f = \delta - \arctan\!\left(\frac{v_y + L_f r}{v_x}\right), \qquad
\alpha_r = -\arctan\!\left(\frac{v_y - L_r r}{v_x}\right)
\]

and cornering stiffnesses $C_f$, $C_r$ identified from vehicle tests.

\subsection{Pacejka ``Magic Formula'' Tire Model}

For cornering beyond the linear regime (slipping or near-limit driving), the
linear tire force model fails.  The \textbf{Pacejka magic formula} (Pacejka,
2012) captures the full nonlinear force--slip curve:

\[
F_y(\alpha) = D\,\sin\!\Big[C\arctan\!\Big(B\alpha - E\big(B\alpha - \arctan(B\alpha)\big)\Big)\Big]
\]

where $B$ (stiffness factor), $C$ (shape factor), $D$ (peak value), and $E$
(curvature factor) are tire-specific parameters.  At small slip angles the
formula reduces to the linear model ($F_y \approx BCD\,\alpha$, so
$C_{\text{eff}} = BCD$), ensuring consistency with the linear model.

Replacing the linear tire forces in the dynamic bicycle model with the Pacejka
formula gives a \textbf{quarter-car + Pacejka} model that is accurate from
low-speed cruising through near-limit cornering.  Both MPPI and the MPC
formulation can accommodate it without structural changes: MPPI requires only
that the forward model be evaluable; MPC requires additionally that $F_y$ is
differentiable, which the smooth Pacejka formula satisfies.


% ─────────────────────────────────────────────────────────────────────────────
\section{Experimental Results}

All three controllers are benchmarked on the same figure-8 raceline (loop
radius $R \approx 1.5$\,m, total arc length $\approx 17.3$\,m, constant
reference speed $v_{\text{ref}} = 4.5$\,m\,s$^{-1}$) over $10$ laps each.
Errors are accumulated per control step from the lap start to the lap end
and reported as lap-averaged RMS values.

\subsection{Metrics}

For each lap $\ell$, with $T_\ell$ samples:
\[
\text{RMS}_e^{(\ell)} = \sqrt{\frac{1}{T_\ell}\sum_{t=1}^{T_\ell} e_t^2},
\qquad
\text{RMS}_{\psi_e}^{(\ell)} = \sqrt{\frac{1}{T_\ell}\sum_{t=1}^{T_\ell} \psi_{e,t}^2},
\qquad
\text{RMS}_{v_e}^{(\ell)} = \sqrt{\frac{1}{T_\ell}\sum_{t=1}^{T_\ell} (v_t - v_{\text{ref},t})^2}.
\]

Aggregate metrics are the means and standard deviations over $\ell = 1, \dots, 10$.

\subsection{Results}

\begin{center}
\begin{tabular}{lcccc}
\toprule
Controller & $|e|_{\text{RMS}}$ (m) & $|\psi_e|_{\text{RMS}}$ (rad) & $|v_e|_{\text{RMS}}$ (m/s) & Notes \\
\midrule
PID  &  --- & --- & --- & Baseline \\
MPPI &  --- & --- & --- & 300 rollouts, $\lambda = 20$ \\
MPC  &  --- & --- & --- & SQP-RTI, ERK4 \\
\bottomrule
\end{tabular}
\end{center}

Relative differences (with PID as the reference):

\begin{center}
\begin{tabular}{lccc}
\toprule
Controller & $\Delta |e|$ vs PID & $\Delta |\psi_e|$ vs PID & $\Delta |v_e|$ vs PID \\
\midrule
MPPI & --- & --- & --- \\
MPC  & --- & --- & --- \\
\bottomrule
\end{tabular}
\end{center}

\subsection{Discussion}

\emph{(To be filled in once the 10-lap runs are collected.)} Expected qualitative
findings, to be confirmed by the data:

\begin{itemize}
  \item \textbf{PID} should have the largest $|e|_{\text{RMS}}$, especially
        through corner entries and exits, because it reacts only to current
        error and cannot anticipate the curvature change.
  \item \textbf{MPPI} should outperform PID by virtue of horizon-aware planning,
        but show non-zero step-to-step variation in $\delta$ because the
        weighted-average update is stochastic.  Steering RMS should be slightly
        higher than MPC.
  \item \textbf{MPC} should produce the smoothest steering trace and the
        smallest tracking errors, at the cost of compile-time complexity and
        a (one-time) solver build step.
  \item All three should track $v_{\text{ref}}$ similarly, since the
        longitudinal channel is governed by the same feedforward-plus-P
        throttle controller in every case.
\end{itemize}

\end{document}
